\setcounter{section}{15}

  \zwischenueberschrift{Übungsaufgaben}

\inputaufgabegibtloesung
{}
{

Es sei $M$ eine kompakte differenzierbare Mannigfaltigkeit mit einer stetigen positiven Volumenform $\omega$. Zeige, dass

\mavergleichskettedisp
{\vergleichskette
{ \int_M \omega
}
{ <} { \infty
}
{ } {
}
{ } {
}
{ } {
}
}
{}{}{}
ist.

}
{} {}

\inputaufgabe
{}
{

Zeige, dass das zu einer
\definitionsverweis {stetigen}{}{}
\definitionsverweis {positiven Volumenform}{}{}
auf einer
\definitionsverweis {differenzierbaren Mannigfaltigkeit}{}{}
in
[[Mannigfaltigkeit/Abzählbar/Positive Volumenform/Zugehöriges Maß/Definition|Definition 15.3]]
eingeführte Volumenmaß ein
$\sigma$-\definitionsverweis {endliches Maß}{}{}
ist.

}
{} {}

\inputaufgabe
{}
{

Es sei

\mavergleichskettedisp
{\vergleichskette
{ \omega
}
{ =} { dx_1 \wedge \ldots \wedge dx_n
}
{ =} { e_1^* \wedge \ldots \wedge e_n^*
}
{ } {
}
{ } {
}
}
{}{}{}
die Standard-Volumenform auf dem $\R^n$. Zeige, dass für jede
\definitionsverweis {messbare Teilmenge}{}{}

\mathl{T \subseteq \R^n}{} die Gleichheit

\mavergleichskettedisp
{\vergleichskette
{
\int_{ T  }   \omega
}
{ =} { \int_{  T  }     \,  d  \lambda^n
}
{ =} { \lambda^n(T)
}
{ } {
}
{ } {
}
}
{}{}{}
gilt.

}
{} {}

\inputaufgabe
{}
{

Es sei $M$ eine
\definitionsverweis {differenzierbare Mannigfaltigkeit}{}{}
mit einer
\definitionsverweis {positiven Volumenform}{}{}
$\omega$. Es sei
\mathl{T \subseteq M}{}
\definitionsverweis {messbar}{}{}
und
\mathl{N \subseteq M}{} eine
\definitionsverweis {Nullmenge}{}{.}
Zeige, dass

\mavergleichskettedisp
{\vergleichskette
{
\int_{ T  }   \omega
}
{ =} {
\int_{ T \setminus N  }   \omega
}
{ } {
}
{ } {
}
{ } {
}
}
{}{}{}
gilt.

}
{} {}

\inputaufgabe
{}
{

Es sei $M$ eine
$n$-\definitionsverweis {dimensionale}{}{}
\definitionsverweis {differenzierbare Mannigfaltigkeit}{}{}
mit einer
\definitionsverweis {abzählbaren Basis der Topologie}{}{}
und es seien
\mathkor {} {\omega_1} {und} {\omega_2} {}
\definitionsverweis {positive Volumenformen}{}{}
auf $M$. Zeige, dass für jede
\definitionsverweis {messbare Teilmenge}{}{}

\mathl{T \subseteq M}{} und
\mathl{a,b \in \R_+}{} die Beziehung

\mavergleichskettedisp
{\vergleichskette
{
\int_{ T  }  (a \omega_1 + b \omega_2)
}
{ =} { a
\int_{  T  }   \omega_1+ b
\int_{  T  }   \omega_2
}
{ } {
}
{ } {
}
{ } {
}
}
{}{}{}
gilt.

}
{} {}

\inputaufgabe
{}
{

Es sei $M$ eine
\definitionsverweis {differenzierbare Mannigfaltigkeit}{}{}
mit
\definitionsverweis {abzählbarer Basis der Topologie}{}{.}
Zeige, wie man unter Bezug auf Karten \anfuehrung{Nullmengen}{} von $M$ erklären kann, ohne dass ein
\definitionsverweis {Maß}{}{}
gegeben ist. Zeige ferner, dass, wenn eine
\definitionsverweis {positive Volumenform}{}{}
gegeben ist, diese Nullmengen auch
\definitionsverweis {Nullmengen}{}{}
im Sinne der Maßtheorie sind.

}
{} {}

\inputaufgabe
{}
{

Beschreibe den Einheitskreis
\mathl{S^1 \subset \R^2}{} als Faser einer Abbildung
\maabbdisp {} {\R^2} {\R
} {}
derart, dass die gemäß
[[Differenzierbare reguläre Abbildung/R^n/Faser besitzt Volumenform über Gradienten/Fakt|Korollar 15.6]]
gegebene Volumenform $\omega$  positiv ist. Berechne
\mathl{\int_{S^1} \omega}{.}

}
{} {}

\inputaufgabe
{}
{

Beschreibe die Einheitssphäre
\mathl{S^2 \subset \R^3}{} als Faser einer Abbildung
\maabbdisp {} {\R^3} {\R
} {}
derart, dass die gemäß
[[Differenzierbare reguläre Abbildung/R^n/Faser besitzt Volumenform über Gradienten/Fakt|Korollar 15.6]]
gegebene Volumenform $\omega$  positiv ist. Berechne
\mathl{\int_{S^2} \omega}{.}

}
{} {}

\inputaufgabe
{}
{

Es seien
\mathkor {} {L} {und} {M} {}
\definitionsverweis {differenzierbare Mannigfaltigkeiten}{}{}
und
\maabbdisp {\varphi} { L } { M
} {}
eine
\definitionsverweis {differenzierbare Abbildung}{}{.}
Es sei

\mavergleichskette
{\vergleichskette
{ \omega
}
{ \in }{
{ \mathcal E }^{  1 } (  M   )
}
{  }{
}
{    }{
}
{   }{
}
}
{}{}{}
eine messbare Differentialform mit der
\definitionsverweis {zurückgezogenen Differentialform}{}{}

\mavergleichskette
{\vergleichskette
{ \varphi^*\omega
}
{ \in }{
{ \mathcal E }^{  1 } (  L   )
}
{  }{
}
{    }{
}
{   }{
}
}
{}{}{}
und es sei
\maabbdisp {\gamma} { I } { L
} {}
eine
\definitionsverweis {stetig differenzierbare Kurve}{}{}
\zusatzklammer {$I$ ein
\definitionsverweis {reelles Intervall}{}{}} {} {.}
Zeige, dass für die
\definitionsverweis {Wegintegrale}{}{}
die Gleichheit

\mavergleichskettedisp
{\vergleichskette
{ \int_\gamma \varphi^* \omega
}
{ =} { \int_{\varphi \circ \gamma} \omega
}
{ } {
}
{ } {
}
{ } {
}
}
{}{}{.}

}
{} {}

\inputaufgabe
{}
{

Sei
\maabbeledisp {\gamma} { [0,2\pi] } {\R^2
} { t } {( \cos  t , \sin  t )
} {,}
gegeben. Berechne das
\definitionsverweis {Wegintegral}{}{}
längs dieses Weges zu den folgenden
\definitionsverweis {Differentialformen}{}{}

a)
\mathl{xdx +ydy}{,}

b)
\mathl{xdx -ydy}{,}

c)
\mathl{ydx +xdy}{,}

d)
\mathl{ydx -xdy}{.}

}
{} {}

\inputaufgabegibtloesung
{}
{

Berechne das Wegintegral
\mathl{\int_\gamma \omega}{} zu
\maabbeledisp {\gamma} { [-1,0] } {\R^3
} { t } {(-t^2,t^3-1,t+2)
} {,}
für die $1$-Differentialform

\mavergleichskettedisp
{\vergleichskette
{ \omega
}
{ =} { x^3dx -yzdy +xz^2 dz
}
{ } {
}
{ } {
}
{ } {
}
}
{}{}{}
auf dem $\R^3$.

}
{} {}

\inputaufgabegibtloesung
{}
{

Wir betrachten die differenzierbaren Abbildungen
\maabbeledisp {\gamma} { [1,2] } {\R^2
} { t } {(t,t^{-1})
} {,}
und
\maabbeledisp {\varphi} {\R^2} {\R^3
} {(u,v)} {(u^2,uv,-u+v^2)
} {,}
und die Differentialform

\mavergleichskettedisp
{\vergleichskette
{ \omega
}
{ =} { xdx-zdy+dz
}
{ } {
}
{ } {
}
{ } {
}
}
{}{}{}
auf dem $\R^3$.

a) Berechne die zurückgezogene Differentialform
\mathl{\varphi^* \omega}{} auf dem $\R^2$.

b) Berechne das Wegintegral zur Differentialform
\mathl{\varphi^* \omega}{} zum Weg $\gamma$.

c) Berechne
\zusatzklammer {ohne Bezug auf b)} {} {}
das Wegintegral zur Differentialform
\mathl{\omega}{} zum Weg $\varphi \circ \gamma$.

}
{} {}

\inputaufgabegibtloesung
{}
{

Wir betrachten die differenzierbaren Abbildungen
\maabbeledisp {\gamma} { [1,c] } {\R^2
} { t } {(t,t^{3})
} {,}
\zusatzklammer {mit \mathlk{c \geq 1}{}} {} {}
und
\maabbeledisp {\varphi} {\R_+ \times \R_+ } {\R^3
} {(u,v)} {(u^3,u^2+v^2,u^{-1}v^{-1} )
} {,}
und die Differentialform

\mavergleichskettedisp
{\vergleichskette
{ \omega
}
{ =} { (x-y)dx-z^2dy+dz
}
{ } {
}
{ } {
}
{ } {
}
}
{}{}{}
auf dem $\R^3$.

a) Berechne die zurückgezogene Differentialform
\mathl{\varphi^* \omega}{} auf dem $\R_+ \times \R_+$.

b) Berechne das Wegintegral zur Differentialform
\mathl{\varphi^* \omega}{} zum Weg $\gamma$ in Abhängigkeit von $c$.

}
{} {}

  \zwischenueberschrift{Aufgaben zum Abgeben}

\inputaufgabe
{4}
{

Es sei $M$ eine
$n$-\definitionsverweis {dimensionale}{}{}
\definitionsverweis {differenzierbare Mannigfaltigkeit}{}{}
mit
\definitionsverweis {abzählbarer Basis der Topologie}{}{.}
Es sei $\omega$ eine
\definitionsverweis {positive Volumenform}{}{} auf $M$ und es sei $\mu$ das durch diese Volumenform definierte
\definitionsverweis {Maß}{}{} auf $M$. Zeige, dass dann jede
\definitionsverweis {abgeschlossene Untermannigfaltigkeit}{}{} der Dimension
\mathl{\leq n-1}{} eine
\definitionsverweis {Nullmenge}{}{}
ist.

}
{} {}

\inputaufgabe
{4}
{

Es seien

\mavergleichskette
{\vergleichskette
{a,b,c,d,r,s
}
{ \geq }{ 1
}
{  }{
}
{    }{
}
{   }{
}
}
{}{}{}
\definitionsverweis {natürliche Zahlen}{}{.}
Wir betrachten die
\definitionsverweis {stetig differenzierbare Kurve}{}{}
\maabbeledisp {} { [0,1] } {\R^2
} { t } {(t^r,t^s)
} {.}
Berechne das
\definitionsverweis {Wegintegral}{}{}
längs dieses Weges zur
\definitionsverweis {Differentialform}{}{}

\mavergleichskette
{\vergleichskette
{ \omega
}
{ = }{ x^ay^b dx +x^c y^d dy
}
{  }{
}
{    }{
}
{   }{
}
}
{}{}{.}

}
{} {}

\inputaufgabe
{5}
{

Es sei
\maabbeledisp {\gamma} { [0,2\pi] } {\R^3
} { t } {( \cos  t , \sin  t, t )
} {,}
gegeben. Berechne das
\definitionsverweis {Wegintegral}{}{}
längs dieses Weges zur
\definitionsverweis {Differentialform}{}{}

\mavergleichskette
{\vergleichskette
{ \omega
}
{ = }{ (y-z^3) dx +x^2dy -xzdz
}
{  }{
}
{    }{
}
{   }{
}
}
{}{}{.}

}
{} {}

 [[Kategorie:Latexseite]]
